% Generated draft: requires mathematical and source review.
\documentclass[11pt]{article}
\usepackage[margin=25mm]{geometry}
\usepackage{fontspec}
\setmainfont{Latin Modern Roman}
\usepackage{amsmath,amssymb,mathtools,unicode-math}
\setmathfont{Latin Modern Math}
\usepackage{enumitem,needspace,xurl,hyperref}
\setlength{\parindent}{0pt}
\setlength{\parskip}{5pt}
\setlength{\emergencystretch}{3em}
\newcommand{\N}{\mathbb{N}}
\newcommand{\Z}{\mathbb{Z}}
\newcommand{\Q}{\mathbb{Q}}
\newcommand{\R}{\mathbb{R}}
\newcommand{\C}{\mathbb{C}}
\newcommand{\F}{\mathbb{F}}
\newcommand{\A}{\mathbb{A}}
\newcommand{\PP}{\mathbb{P}}
\newcommand{\E}{\mathbb{E}}
\newcommand{\eps}{\varepsilon}
\newcommand{\dd}{\,\mathrm{d}}
\begin{document}
\section*{UnsolvedMath ID: 30006463}
\textbf{Source problem code:} catalogue entry
\subsection{Definitions and mathematical statement}
Let \(a=(a(n))_{n\geq 1}\) and \(b=(b(n))_{n\geq 1}\) denote Fourier coefficient sequences of automorphic forms in the intended, but unspecified, higher-rank setting. Using positive-integer indexing, define the shift-one convolution sum by
\[
S_{a,b}(x):=\sum_{\substack{n\in\Z\\1\leq n\leq x}}a(n)b(n+1),\qquad x\in\R,\quad x\geq 1.
\]
Positive-integer indexing is an explicit interpretation of the displayed source formula; the extraction does not state the lower summation boundary. The product is written without complex conjugation, as supplied.

\emph{Research target.} Obtain nontrivial estimates and, in appropriate cases, asymptotic formulas for \(S_{a,b}(x)\) as \(x\to\infty\), extending the cases already understood for automorphic Fourier coefficients in lower rank.

\emph{Unresolved scope.} The record does not specify the automorphic groups or ranks, whether one or both factors must have higher rank, the class of forms, their coefficient indexing and normalization, or any uniformity in form parameters. It also supplies neither a comparison bound defining ``nontrivial'' nor a proposed main term, error scale, or quantitative saving. Consequently, this is a research target rather than a fully quantified conjecture. No assertion for every higher-rank pair, and no particular asymptotic formula, can be extracted reliably.

The original discussion begins with general arithmetic functions, usually multiplicative, and mentions connections with the twin prime and Chowla conjectures, moments of \(L\)-functions, subconvexity, and quantum unique ergodicity. These are contextual motivations, not additional hypotheses or claims of equivalence. The supplied dated assessment reports progress in important \(\mathrm{GL}(3)\times\mathrm{GL}(2)\) cases, including a remaining divisor-function case, while labeling the general extracted target open. These are supplied literature qualifications and have not been independently verified.
\subsection{Short English statement}
Study cancellation and asymptotic behavior in the shift-one sums of Fourier coefficients of higher-rank automorphic forms. The supplied record labels the general target open, but leaves the classes of forms and the desired bounds and main terms unspecified; a precise quantified conjecture therefore remains unresolved.
\subsection{Sources}
\begin{itemize}
\item Supplied local UnsolvedMath registry record, permanent ID \(30006463\), title ``Higher-Rank Shifted Convolution Sums,'' source problem code \texttt{catalogue entry}, in the Oberwolfach Reports Open Problems collection. The local registry locations are \texttt{data/unsolvedmath/problems.json} and the supplementary \texttt{catalog/problems.json}; this draft uses the supplied complete record without independently checking those files.
\item Supplied report citation: \emph{Analytic Number Theory (2026)}, Oberwolfach Reports, DOI \texttt{10.4171/owr/2025/51}, pp.\ 2701--2764. Supplied source URL: \url{https://doi.org/10.4171/owr/2025/51}. The report title's year and the DOI's year component differ in the supplied record and have not been reconciled. No individual proposer is identified; the registry gives proposed year \(2025\).
\item Supplied literature source: ``Oberwolfach report 51/2025,'' \url{https://doi.org/10.4171/owr/2025/51}. The exact supplied literature-check date is \texttt{2026-08-22}. This is supplied dated curation, not an independently verified assessment. It labels the general higher-rank formulation open and reports known asymptotics in important \(\mathrm{GL}(3)\times\mathrm{GL}(2)\) cases, including a remaining divisor-function case. The record expressly requests independent verification of time-sensitive status and cited claims.
\item The registry labels the statement \texttt{corrected\_verified} and describes repairs to formula delimiters, notation, grammar, and typography while preserving mathematical scope. That supplied editorial label does not establish mathematical correctness or independent source verification.
\end{itemize}
\end{document}
